\begin{proofbox}
	\textbf{\textcolor{CProof}{证明一：从对数的定义出发（最简单、最重要）}}
	已知 $a>0$、$a\neq 1$、$N>0$。
\begin{center}
\begin{tikzpicture}[font=\small,>=stealth]
\draw[draw=teal!50!black,fill=green!3,rounded corners=3pt] (0,0) rectangle (13,6);
\node at (6.5,5.45) {同一条指数曲线：已知高度找指数，已知指数找高度};
\begin{scope}[shift={(1.3,1.2)},x=1.1cm,y=.4cm]
\draw[gray,->] (-.6,0)--(3.8,0) node[right] {$x$};\draw[gray,->] (0,0)--(0,9.5) node[above] {$y$};
\draw[teal!70!black,very thick,domain=-.5:3.2,samples=70] plot(\x,{2^\x});
\draw[->,blue!70!black,thick] (0,8)--(2.95,8);
\draw[->,orange!85!black,thick] (3,7.9)--(3,.15);
\fill[teal!70!black] (3,8) circle(2.7pt);
\node[above left] at (3,8.3) {$(3,8)$};\node[left,blue!70!black] at (0,8) {$N=8$};
\node[below,orange!85!black] at (3,-.05) {$t=3$};\node[left] at (0,1) {$1$};
\end{scope}
\node[anchor=west,align=left,text width=5.4cm] at (6.8,3.1) {$y=2^x$ 上只有一个点的高度是 $8$。\\[10pt]横着读：$\log_2 8=3$。\\[10pt]竖着读：$2^3=8$。\\[10pt]两次读取的是同一个点。};
\node at (6.5,.45) {先找指数再算高度，或先算高度再找指数，都回到原数。};
\end{tikzpicture}

\end{center}
\textbf{读图：}从高度 $N$ 水平找到曲线，再向下读取指数 $t$；从这个 $t$ 又能读回同一高度 $N$。
	\begin{description}[style=nextline, leftmargin=2.8em, labelsep=0.5em, itemsep=0.4em]
		\item[\textbf{第一步：给未知的指数起个名字。}]
			设 $t=\log_a N$。这句话的意思是：“$a$ 的 $t$ 次方等于 $N$”。
		\item[\textbf{第二步：按定义写成乘方。}]
			因此 $a^t=N$。
		\item[\textbf{第三步：把 $t$ 换回原样。}]
			因为 $t=\log_a N$，于是
			\[
			\boxed{a^{\log_a N}=a^t=N}.
			\]
	\end{description}
	\textbf{为什么这就够了？}因为 $\log_a N$ 本来就是“使 $a$ 的乘方等于 $N$ 的那个指数”。这个等式不是新规定，而是对数定义的直接结果。

	\medskip
	\textbf{\textcolor{CProof}{证明二：把顺序反过来}}
\begin{center}
\begin{tikzpicture}[font=\small,>=stealth]
\draw[draw=teal!50!black,fill=green!3,rounded corners=3pt] (0,0) rectangle (13,7.2);
\node at (6.5,6.65) {交换一个点的横、纵坐标：指数与对数把输入、输出互换};
\begin{scope}[shift={(1.2,1.2)},x=.55cm,y=.55cm]
\draw[gray,->] (-.4,0)--(9.5,0) node[right] {$x$};\draw[gray,->] (0,-.4)--(0,9.5) node[above] {$y$};
\draw[gray,dashed] (0,0)--(9,9);\node[below right] at (7.9,8.2) {$y=x$};
\draw[teal!70!black,very thick,domain=0:3.17,samples=70] plot(\x,{2^\x});
\draw[orange!85!black,very thick,domain=1:9,samples=80] plot(\x,{ln(\x)/ln(2)});
\draw[blue!60!black,dotted] (3,8)--(8,3);
\fill[teal!70!black] (3,8) circle(2.6pt);\fill[orange!85!black] (8,3) circle(2.6pt);
\node[above left] at (3,8.2) {$(3,8)$};\node[below] at (8,2.8) {$(8,3)$};
\end{scope}
\node[anchor=west,align=left,text width=5.1cm] at (7.3,3.8) {绿线 $y=2^x$：输入 $3$，输出 $8$。\\[12pt]橙线 $y=\log_2x$：输入 $8$，输出 $3$。\\[12pt]因此 $\log_2(2^3)=3$。\\[12pt]一般底数 $a>0,\ a\ne1$ 同理。};
\node at (6.5,.45) {指数函数的值总是正数，正好能作为对数的输入。};
\end{tikzpicture}

\end{center}
	设 $N=a^x$。因 $a>0$，所以 $a^x>0$，这个对数一定有意义。
	\[
	\log_a(a^x)=\log_a N=x.
	\]
	\textbf{理由：}既然 $a^x=N$，那么“$a$ 的几次方等于 $N$”的答案就是 $x$。这说明第二个恒等式对\textbf{任意实数 $x$} 都成立。

	\medskip
	\textbf{\textcolor{CProof}{再用具体数字看一遍（帮助理解，不代替证明）}}
	\begin{center}
	\begin{tikzpicture}[>=stealth, font=\small]
		\node[draw=CProof, rounded corners=3pt, minimum width=2.0cm, minimum height=0.8cm] (u) at (0,0) {$2^3=8$};
		\node[draw=CProof, rounded corners=3pt, minimum width=2.2cm, minimum height=0.8cm] (v) at (5.8,0) {$\log_2 8=3$};
		\draw[->, thick, draw=CProof, bend left=15] (u.north east) to node[above] {已知乘方，反过来找指数} (v.north west);
		\draw[->, thick, draw=CProof, bend left=15] (v.south west) to node[below] {已知指数，算回原来的数} (u.south east);
	\end{tikzpicture}

	\end{center}
\end{proofbox}
